% --- Archon physics formalization source begin ---
% archon:physics
% archon:covers PhyXMiniProblems/problem_phyx_mini_0271.lean
% archon:source-report reports/phyx_mini/problem_phyx_mini_0271.source.json

\chapter{Physics problem phyx\_mini\_0271}
\label{ch:PhyXMiniProblems_problem_phyx_mini_0271}

\paragraph{Problem source.}
\#\# Physical scenario

In figure, three \$10000\textbackslash{} kg\$ ore cars are held at rest on a mine railway using a cable that is parallel to the rails, which are inclined at angle \$\textbackslash{}theta = 30\textasciicircum{}\textbackslash{}circ\$. The cable stretches \$15\textbackslash{} cm\$ just before the coupling between the two lower cars breaks, detaching the lowest car.

\#\# Question

Assuming that the cable obeys Hooke's law, find the amplitude of the resulting oscillations of the remaining two cars.

\#\# Answer choices

- A: A: 0.030 m

- B: B: 0.040 m

- C: C: 0.050 m

- D: D: 0.060 m

\#\# Figure caption (auxiliary; use the image as primary evidence)

The image depicts an inclined surface with three mine carts on a track, loaded with rocks. The track is aligned with a pulley system at the top right. The angle of inclination is labeled as \textbackslash{}( \textbackslash{}theta \textbackslash{}) at the bottom left. The second cart (from the left) is labeled as "Car that breaks free." The area outside the track shows a textured surface resembling dirt or gravel.

\#\# Dataset metadata

Wave Properties; Physical Model Grounding Reasoning, Multi-Formula Reasoning

\paragraph{Recorded answer/context.}
C: C: 0.050 m

\paragraph{Figure/image path.}
/root/proposal\_for\_physic/science-mango/phyx\_mini\_run/phyx\_data/test\_image/271.png

\paragraph{Formalization target.}
create a compiling Lean file with sorry bodies at `PhyXMiniProblems/problem\_phyx\_mini\_0271.lean`. The Lean declarations must preserve the physical quantities, dimensions or dimensional roles, figure labels, governing-law hypotheses, and final relation expressed by this problem.
Use Mathlib/Physlib names found through LeanExplore where available. If a physics API is missing, introduce faithful local abstractions rather than scalar placeholder aliases.

\begin{theorem}[Physics formalization target]
\label{thm:physics:phyx_mini_0271:target}
\lean{PhyXMiniProblems.ProblemPhyXMini0271.oscillationAmplitude_is_five_centimeters_and_matches_answerC}
\uses{def:physics:phyx-mini-0271:phyxminiproblems-problemphyxmini0271-lengthinmeters, def:physics:phyx-mini-0271:phyxminiproblems-problemphyxmini0271-minecarcablesetup, def:physics:phyx-mini-0271:phyxminiproblems-problemphyxmini0271-matchesproblemdata, def:physics:phyx-mini-0271:phyxminiproblems-problemphyxmini0271-matchessuppliedfigure, def:physics:phyx-mini-0271:phyxminiproblems-problemphyxmini0271-hasphysicalparameters, def:physics:phyx-mini-0271:phyxminiproblems-problemphyxmini0271-satisfieselasticcableoscillatorlaws, def:physics:phyx-mini-0271:phyxminiproblems-problemphyxmini0271-recordedanswerchoice, def:physics:phyx-mini-0271:phyxminiproblems-problemphyxmini0271-matchesanswerchoice}
The assigned autoformalize agent should translate this physics problem into Lean declarations in the covered file, with theorem and lemma proof bodies written as `by sorry`.
\end{theorem}
\begin{proof}
This is an autoformalization task, not a proof task. Produce statements that can later be proved by the physics prover without weakening the physical model.
\end{proof}

% --- Archon declaration topology begin ---
\section{Declaration topology}
\begin{definition}[Mass Quantity]
  \label{def:physics:phyx-mini-0271:phyxminiproblems-problemphyxmini0271-massquantity}
  \lean{PhyXMiniProblems.ProblemPhyXMini0271.MassQuantity}
  A nonnegative physical mass
\end{definition}
\begin{proof}
  This is the declared model definition of Mass Quantity.
\end{proof}
\begin{definition}[Length Quantity]
  \label{def:physics:phyx-mini-0271:phyxminiproblems-problemphyxmini0271-lengthquantity}
  \lean{PhyXMiniProblems.ProblemPhyXMini0271.LengthQuantity}
  A nonnegative physical length, used for cable extensions and amplitude
\end{definition}
\begin{proof}
  This is the declared model definition of Length Quantity.
\end{proof}
\begin{definition}[Force Quantity]
  \label{def:physics:phyx-mini-0271:phyxminiproblems-problemphyxmini0271-forcequantity}
  \lean{PhyXMiniProblems.ProblemPhyXMini0271.ForceQuantity}
  A nonnegative physical force
\end{definition}
\begin{proof}
  This is the declared model definition of Force Quantity.
\end{proof}
\begin{definition}[Acceleration Quantity]
  \label{def:physics:phyx-mini-0271:phyxminiproblems-problemphyxmini0271-accelerationquantity}
  \lean{PhyXMiniProblems.ProblemPhyXMini0271.AccelerationQuantity}
  A nonnegative acceleration
\end{definition}
\begin{proof}
  This is the declared model definition of Acceleration Quantity.
\end{proof}
\begin{definition}[Speed Quantity]
  \label{def:physics:phyx-mini-0271:phyxminiproblems-problemphyxmini0271-speedquantity}
  \lean{PhyXMiniProblems.ProblemPhyXMini0271.SpeedQuantity}
  A nonnegative speed
\end{definition}
\begin{proof}
  This is the declared model definition of Speed Quantity.
\end{proof}
\begin{definition}[Cable Stiffness Quantity]
  \label{def:physics:phyx-mini-0271:phyxminiproblems-problemphyxmini0271-cablestiffnessquantity}
  \lean{PhyXMiniProblems.ProblemPhyXMini0271.CableStiffnessQuantity}
  A cable stiffness, with dimension force per length
\end{definition}
\begin{proof}
  This is the declared model definition of Cable Stiffness Quantity.
\end{proof}
\begin{definition}[mass In Kilograms]
  \label{def:physics:phyx-mini-0271:phyxminiproblems-problemphyxmini0271-massinkilograms}
  \lean{PhyXMiniProblems.ProblemPhyXMini0271.massInKilograms}
  \uses{def:physics:phyx-mini-0271:phyxminiproblems-problemphyxmini0271-massquantity}
  Kilogram readout of a physical mass
\end{definition}
\begin{proof}
  This is the declared model definition of mass In Kilograms.
\end{proof}
\begin{definition}[length In Meters]
  \label{def:physics:phyx-mini-0271:phyxminiproblems-problemphyxmini0271-lengthinmeters}
  \lean{PhyXMiniProblems.ProblemPhyXMini0271.lengthInMeters}
  \uses{def:physics:phyx-mini-0271:phyxminiproblems-problemphyxmini0271-lengthquantity}
  Metre readout of a physical length
\end{definition}
\begin{proof}
  This is the declared model definition of length In Meters.
\end{proof}
\begin{definition}[force In Newtons]
  \label{def:physics:phyx-mini-0271:phyxminiproblems-problemphyxmini0271-forceinnewtons}
  \lean{PhyXMiniProblems.ProblemPhyXMini0271.forceInNewtons}
  \uses{def:physics:phyx-mini-0271:phyxminiproblems-problemphyxmini0271-forcequantity}
  Newton readout of a physical force
\end{definition}
\begin{proof}
  This is the declared model definition of force In Newtons.
\end{proof}
\begin{definition}[acceleration In Meters Per Second Squared]
  \label{def:physics:phyx-mini-0271:phyxminiproblems-problemphyxmini0271-accelerationinmeterspersecondsquared}
  \lean{PhyXMiniProblems.ProblemPhyXMini0271.accelerationInMetersPerSecondSquared}
  \uses{def:physics:phyx-mini-0271:phyxminiproblems-problemphyxmini0271-accelerationquantity}
  Metres-per-second-squared readout of a physical acceleration
\end{definition}
\begin{proof}
  This is the declared model definition of acceleration In Meters Per Second Squared.
\end{proof}
\begin{definition}[speed In Meters Per Second]
  \label{def:physics:phyx-mini-0271:phyxminiproblems-problemphyxmini0271-speedinmeterspersecond}
  \lean{PhyXMiniProblems.ProblemPhyXMini0271.speedInMetersPerSecond}
  \uses{def:physics:phyx-mini-0271:phyxminiproblems-problemphyxmini0271-speedquantity}
  Metres-per-second readout of a physical speed
\end{definition}
\begin{proof}
  This is the declared model definition of speed In Meters Per Second.
\end{proof}
\begin{definition}[cable Stiffness In Newtons Per Meter]
  \label{def:physics:phyx-mini-0271:phyxminiproblems-problemphyxmini0271-cablestiffnessinnewtonspermeter}
  \lean{PhyXMiniProblems.ProblemPhyXMini0271.cableStiffnessInNewtonsPerMeter}
  \uses{def:physics:phyx-mini-0271:phyxminiproblems-problemphyxmini0271-cablestiffnessquantity}
  Newton-per-metre readout of the elastic cable's stiffness
\end{definition}
\begin{proof}
  This is the declared model definition of cable Stiffness In Newtons Per Meter.
\end{proof}
\begin{definition}[degrees To Radians]
  \label{def:physics:phyx-mini-0271:phyxminiproblems-problemphyxmini0271-degreestoradians}
  \lean{PhyXMiniProblems.ProblemPhyXMini0271.degreesToRadians}
  Convert an angle stated in degrees to its dimensionless radian value
\end{definition}
\begin{proof}
  This is the declared model definition of degrees To Radians.
\end{proof}
\begin{definition}[Ore Car]
  \label{def:physics:phyx-mini-0271:phyxminiproblems-problemphyxmini0271-orecar}
  \lean{PhyXMiniProblems.ProblemPhyXMini0271.OreCar}
  The three ore cars, ordered by position along the incline
\end{definition}
\begin{proof}
  This is the declared model definition of Ore Car.
\end{proof}
\begin{definition}[Load Stage]
  \label{def:physics:phyx-mini-0271:phyxminiproblems-problemphyxmini0271-loadstage}
  \lean{PhyXMiniProblems.ProblemPhyXMini0271.LoadStage}
  The load configurations relevant to the cable equilibrium
\end{definition}
\begin{proof}
  This is the declared model definition of Load Stage.
\end{proof}
\begin{definition}[Car Coupling]
  \label{def:physics:phyx-mini-0271:phyxminiproblems-problemphyxmini0271-carcoupling}
  \lean{PhyXMiniProblems.ProblemPhyXMini0271.CarCoupling}
  The two couplings between adjacent cars
\end{definition}
\begin{proof}
  This is the declared model definition of Car Coupling.
\end{proof}
\begin{definition}[Cable Orientation]
  \label{def:physics:phyx-mini-0271:phyxminiproblems-problemphyxmini0271-cableorientation}
  \lean{PhyXMiniProblems.ProblemPhyXMini0271.CableOrientation}
  Orientation of the cable relative to the mine railway
\end{definition}
\begin{proof}
  This is the declared model definition of Cable Orientation.
\end{proof}
\begin{definition}[Oscillation Model]
  \label{def:physics:phyx-mini-0271:phyxminiproblems-problemphyxmini0271-oscillationmodel}
  \lean{PhyXMiniProblems.ProblemPhyXMini0271.OscillationModel}
  Idealized model for the motion after the coupling failure
\end{definition}
\begin{proof}
  This is the declared model definition of Oscillation Model.
\end{proof}
\begin{definition}[Supplied Figure]
  \label{def:physics:phyx-mini-0271:phyxminiproblems-problemphyxmini0271-suppliedfigure}
  \lean{PhyXMiniProblems.ProblemPhyXMini0271.SuppliedFigure}
  \uses{def:physics:phyx-mini-0271:phyxminiproblems-problemphyxmini0271-orecar}
  Qualitative components and annotations visible in the supplied figure
\end{definition}
\begin{proof}
  This is the declared model definition of Supplied Figure.
\end{proof}
\begin{definition}[Mine Car Cable Setup]
  \label{def:physics:phyx-mini-0271:phyxminiproblems-problemphyxmini0271-minecarcablesetup}
  \lean{PhyXMiniProblems.ProblemPhyXMini0271.MineCarCableSetup}
  \uses{def:physics:phyx-mini-0271:phyxminiproblems-problemphyxmini0271-massquantity, def:physics:phyx-mini-0271:phyxminiproblems-problemphyxmini0271-lengthquantity, def:physics:phyx-mini-0271:phyxminiproblems-problemphyxmini0271-forcequantity, def:physics:phyx-mini-0271:phyxminiproblems-problemphyxmini0271-accelerationquantity, def:physics:phyx-mini-0271:phyxminiproblems-problemphyxmini0271-speedquantity, def:physics:phyx-mini-0271:phyxminiproblems-problemphyxmini0271-cablestiffnessquantity, def:physics:phyx-mini-0271:phyxminiproblems-problemphyxmini0271-orecar, def:physics:phyx-mini-0271:phyxminiproblems-problemphyxmini0271-loadstage, def:physics:phyx-mini-0271:phyxminiproblems-problemphyxmini0271-carcoupling, def:physics:phyx-mini-0271:phyxminiproblems-problemphyxmini0271-cableorientation, def:physics:phyx-mini-0271:phyxminiproblems-problemphyxmini0271-oscillationmodel, def:physics:phyx-mini-0271:phyxminiproblems-problemphyxmini0271-suppliedfigure}
  The complete apparatus and the unknown post-break quantities. The new equilibrium extension and the oscillation amplitude are independent fields. In particular, neither is defined from the recorded answer
\end{definition}
\begin{proof}
  This is the declared model definition of Mine Car Cable Setup.
\end{proof}
\begin{definition}[attached Mass In Kilograms]
  \label{def:physics:phyx-mini-0271:phyxminiproblems-problemphyxmini0271-attachedmassinkilograms}
  \lean{PhyXMiniProblems.ProblemPhyXMini0271.attachedMassInKilograms}
  \uses{def:physics:phyx-mini-0271:phyxminiproblems-problemphyxmini0271-massinkilograms, def:physics:phyx-mini-0271:phyxminiproblems-problemphyxmini0271-orecar, def:physics:phyx-mini-0271:phyxminiproblems-problemphyxmini0271-loadstage, def:physics:phyx-mini-0271:phyxminiproblems-problemphyxmini0271-minecarcablesetup}
  Total mass carried by the cable at a given stage, read in kilograms
\end{definition}
\begin{proof}
  This is the declared model definition of attached Mass In Kilograms.
\end{proof}
\begin{definition}[Matches Problem Data]
  \label{def:physics:phyx-mini-0271:phyxminiproblems-problemphyxmini0271-matchesproblemdata}
  \lean{PhyXMiniProblems.ProblemPhyXMini0271.MatchesProblemData}
  \uses{def:physics:phyx-mini-0271:phyxminiproblems-problemphyxmini0271-massinkilograms, def:physics:phyx-mini-0271:phyxminiproblems-problemphyxmini0271-lengthinmeters, def:physics:phyx-mini-0271:phyxminiproblems-problemphyxmini0271-speedinmeterspersecond, def:physics:phyx-mini-0271:phyxminiproblems-problemphyxmini0271-minecarcablesetup}
  Numerical and qualitative facts stated by the problem. The gravity and cable stiffness are deliberately not assigned numerical values because both cancel when the two static equilibria are compared
\end{definition}
\begin{proof}
  This is the declared model definition of Matches Problem Data.
\end{proof}
\begin{definition}[Matches Supplied Figure]
  \label{def:physics:phyx-mini-0271:phyxminiproblems-problemphyxmini0271-matchessuppliedfigure}
  \lean{PhyXMiniProblems.ProblemPhyXMini0271.MatchesSuppliedFigure}
  \uses{def:physics:phyx-mini-0271:phyxminiproblems-problemphyxmini0271-minecarcablesetup}
  Primary-image readout: three cars lie on an inclined rail, the cable from the upper car runs along the rails toward a pulley, θ marks the incline, and the break-free annotation identifies the lowest car
\end{definition}
\begin{proof}
  This is the declared model definition of Matches Supplied Figure.
\end{proof}
\begin{definition}[Has Physical Parameters]
  \label{def:physics:phyx-mini-0271:phyxminiproblems-problemphyxmini0271-hasphysicalparameters}
  \lean{PhyXMiniProblems.ProblemPhyXMini0271.HasPhysicalParameters}
  \uses{def:physics:phyx-mini-0271:phyxminiproblems-problemphyxmini0271-massinkilograms, def:physics:phyx-mini-0271:phyxminiproblems-problemphyxmini0271-accelerationinmeterspersecondsquared, def:physics:phyx-mini-0271:phyxminiproblems-problemphyxmini0271-cablestiffnessinnewtonspermeter, def:physics:phyx-mini-0271:phyxminiproblems-problemphyxmini0271-minecarcablesetup}
  Positivity and nondegeneracy of the physical inputs
\end{definition}
\begin{proof}
  This is the declared model definition of Has Physical Parameters.
\end{proof}
\begin{definition}[Satisfies Elastic Cable Oscillator Laws]
  \label{def:physics:phyx-mini-0271:phyxminiproblems-problemphyxmini0271-satisfieselasticcableoscillatorlaws}
  \lean{PhyXMiniProblems.ProblemPhyXMini0271.SatisfiesElasticCableOscillatorLaws}
  \uses{def:physics:phyx-mini-0271:phyxminiproblems-problemphyxmini0271-lengthinmeters, def:physics:phyx-mini-0271:phyxminiproblems-problemphyxmini0271-forceinnewtons, def:physics:phyx-mini-0271:phyxminiproblems-problemphyxmini0271-accelerationinmeterspersecondsquared, def:physics:phyx-mini-0271:phyxminiproblems-problemphyxmini0271-speedinmeterspersecond, def:physics:phyx-mini-0271:phyxminiproblems-problemphyxmini0271-cablestiffnessinnewtonspermeter, def:physics:phyx-mini-0271:phyxminiproblems-problemphyxmini0271-degreestoradians, def:physics:phyx-mini-0271:phyxminiproblems-problemphyxmini0271-minecarcablesetup, def:physics:phyx-mini-0271:phyxminiproblems-problemphyxmini0271-attachedmassinkilograms}
  Hooke's law and force balance for the old and new static equilibria, followed by the general Physlib harmonic-oscillator model for release about the new equilibrium. The amplitude field is linked only to Physlib's general amplitude computed from initial position and velocity. No numerical post-break extension, amplitude, or answer choice occurs in these laws
\end{definition}
\begin{proof}
  This is the declared model definition of Satisfies Elastic Cable Oscillator Laws.
\end{proof}
\begin{lemma}[post Break Equilibrium Extension is ten centimeters]
  \label{lem:physics:phyx-mini-0271:phyxminiproblems-problemphyxmini0271-postbreakequilibriumextension-is-ten-centimeters}
  \lean{PhyXMiniProblems.ProblemPhyXMini0271.postBreakEquilibriumExtension_is_ten_centimeters}
  \uses{def:physics:phyx-mini-0271:phyxminiproblems-problemphyxmini0271-lengthinmeters, def:physics:phyx-mini-0271:phyxminiproblems-problemphyxmini0271-minecarcablesetup, def:physics:phyx-mini-0271:phyxminiproblems-problemphyxmini0271-matchesproblemdata, def:physics:phyx-mini-0271:phyxminiproblems-problemphyxmini0271-hasphysicalparameters, def:physics:phyx-mini-0271:phyxminiproblems-problemphyxmini0271-satisfieselasticcableoscillatorlaws}
  Comparing Hooke equilibrium for three equal cars with that for the remaining two cars gives a new cable extension of 10 cm
\end{lemma}
\begin{proof}
  The conclusion follows from the governing relations and figure readouts named in the statement of post Break Equilibrium Extension is ten centimeters.
\end{proof}
\begin{lemma}[oscillation Amplitude eq equilibrium Shift]
  \label{lem:physics:phyx-mini-0271:phyxminiproblems-problemphyxmini0271-oscillationamplitude-eq-equilibriumshift}
  \lean{PhyXMiniProblems.ProblemPhyXMini0271.oscillationAmplitude_eq_equilibriumShift}
  \uses{def:physics:phyx-mini-0271:phyxminiproblems-problemphyxmini0271-lengthinmeters, def:physics:phyx-mini-0271:phyxminiproblems-problemphyxmini0271-minecarcablesetup, def:physics:phyx-mini-0271:phyxminiproblems-problemphyxmini0271-matchesproblemdata, def:physics:phyx-mini-0271:phyxminiproblems-problemphyxmini0271-satisfieselasticcableoscillatorlaws}
  For an undamped linear oscillator released from rest, Physlib's amplitude--phase conversion makes the amplitude the magnitude of the initial displacement from the new equilibrium
\end{lemma}
\begin{proof}
  The conclusion follows from the governing relations and figure readouts named in the statement of oscillation Amplitude eq equilibrium Shift.
\end{proof}
\begin{definition}[Answer Choice]
  \label{def:physics:phyx-mini-0271:phyxminiproblems-problemphyxmini0271-answerchoice}
  \lean{PhyXMiniProblems.ProblemPhyXMini0271.AnswerChoice}
  Labels of the four displayed answers
\end{definition}
\begin{proof}
  This is the declared model definition of Answer Choice.
\end{proof}
\begin{definition}[meters]
  \label{def:physics:phyx-mini-0271:phyxminiproblems-problemphyxmini0271-answerchoice-meters}
  \lean{PhyXMiniProblems.ProblemPhyXMini0271.AnswerChoice.meters}
  \uses{def:physics:phyx-mini-0271:phyxminiproblems-problemphyxmini0271-answerchoice}
  Oscillation amplitude in metres printed beside each answer label
\end{definition}
\begin{proof}
  This is the declared model definition of meters.
\end{proof}
\begin{definition}[recorded Answer Choice]
  \label{def:physics:phyx-mini-0271:phyxminiproblems-problemphyxmini0271-recordedanswerchoice}
  \lean{PhyXMiniProblems.ProblemPhyXMini0271.recordedAnswerChoice}
  \uses{def:physics:phyx-mini-0271:phyxminiproblems-problemphyxmini0271-answerchoice}
  The answer label recorded in the supplied dataset metadata
\end{definition}
\begin{proof}
  This is the declared model definition of recorded Answer Choice.
\end{proof}
\begin{definition}[Matches Answer Choice]
  \label{def:physics:phyx-mini-0271:phyxminiproblems-problemphyxmini0271-matchesanswerchoice}
  \lean{PhyXMiniProblems.ProblemPhyXMini0271.MatchesAnswerChoice}
  \uses{def:physics:phyx-mini-0271:phyxminiproblems-problemphyxmini0271-lengthquantity, def:physics:phyx-mini-0271:phyxminiproblems-problemphyxmini0271-lengthinmeters, def:physics:phyx-mini-0271:phyxminiproblems-problemphyxmini0271-answerchoice}
  Exact agreement of a physical amplitude with a displayed answer
\end{definition}
\begin{proof}
  This is the declared model definition of Matches Answer Choice.
\end{proof}
% --- Archon declaration topology end ---
% --- Archon physics formalization source end ---
